\section{Signed multilinear forms and conductor decay}
\label{sec:functional}

The tuple laws provide two estimates for functions on products of prime
fields. The first follows the reflection and chessboard method
of~\cite[Theorem~2.2]{FrohlichLieb1978} and
\cite[Theorem~4.1]{FILS1978ReflectionI}; the finite proof below verifies
the required inequality for this particular tuple law.  Reflection positivity gives a H\"older inequality for signed
inputs and a seminorm that dominates the uniform mean.  The conditional
mixing estimate gives a separate bound for inputs supported on large
Fourier denominators.  This second estimate also applies to the
nonexceptional laws.

Retain the index set \(V=\Pi^\ell\), where \(\Pi\) consists of the
permutation words on \([h]\), \(h=24\), and \(\ell=161\), and put
\(d=|V|\).  In particular, \(d\) is even.  Fix the prime threshold
\(P\) from Proposition~\ref{prop:tuple-laws} and
Lemma~\ref{lem:conditional-mixing}, and put \(\gamma=1/64\).
For a finite set \(\mathcal P\) of primes at least \(P\), define
\[
 X_{\mathcal P}=\prod_{p\in\mathcal P}\mathbb F_p,
 \qquad
 \mu_{\mathcal P}=\bigotimes_{p\in\mathcal P}\mu_p.
\]
We regard \(\mu_{\mathcal P}\) as a probability law on
\(X_{\mathcal P}^V\) by collecting, for each slot \(v\), its labels
at all primes into \(z_v\in X_{\mathcal P}\).  Every single-slot
marginal is uniform on \(X_{\mathcal P}\).  Unless another measure
is specified, norms and expectations on \(X_{\mathcal P}\) use
uniform probability measure.  Set
\begin{equation}
 \mathcal Z_{\mathcal P}((g_v)_{v\in V})
   =\int\prod_{v\in V}g_v(z_v)\,d\mu_{\mathcal P}(z),
 \qquad
 Z_{\mathcal P}(g)=\mathcal Z_{\mathcal P}((g)_{v\in V}).
 \label{eq:functional-definition}
\end{equation}
For an empty prime set, \(X_{\varnothing}\) is a singleton and
the product law has mass one; thus \(Z_{\varnothing}(g)=g^d\).

\subsection{A finite signed reflection inequality}

\begin{proposition}[Signed H\"older inequality]
\label{prop:signed-holder}
For every finite prime set \(\mathcal P\) as above and real functions
\(g,g_v:X_{\mathcal P}\to\mathbb R\), one has
\begin{equation}
 Z_{\mathcal P}(g)\ge0,
 \qquad
 \left|\mathcal Z_{\mathcal P}((g_v)_{v\in V})\right|
 \le\prod_{v\in V}Z_{\mathcal P}(g_v)^{1/d}.
 \label{eq:signed-holder}
\end{equation}
The inequality includes inputs with zero diagonal integral.
\end{proposition}

\begin{proof}
Across any reflection cut, the matrix of \(\mu_{\mathcal P}\) is
the tensor product of the corresponding matrices of the \(\mu_p\).
It is therefore symmetric positive semidefinite.  This also holds for
the empty prime set.  We first record precisely what a reflection does
to an assignment of functions.

Fix a block and two distinct symbols.  Let \(\theta\) exchange those
symbols in that block, and let \(V_+\) be the vertices in which the
first symbol occurs before the second.  Put \(V_-=\theta V_+\).
Define the two fold maps \(\phi_+,\phi_-:V\to V\) by
\[
 \phi_+(v)=
 \begin{cases}v,&v\in V_+,\\\theta v,&v\in V_-,\end{cases}
 \qquad
 \phi_-(v)=
 \begin{cases}\theta v,&v\in V_+,\\v,&v\in V_-.\end{cases}
\]
For an assignment \(a:V\to\mathbb R^{X_{\mathcal P}}\), abbreviate
\(\mathcal Z_{\mathcal P}((a(v))_{v\in V})\) by \(\Lambda(a)\).
The bilinear form of the cut matrix is evaluated on the two real
functions
\[
 F(x)=\prod_{v\in V_+}a(v)(x_v),
 \qquad
 G(x)=\prod_{v\in V_+}a(\theta v)(x_v).
\]
Its mixed value is \(\Lambda(a)\); its two quadratic values are
\(\Lambda(a\circ\phi_+)\) and \(\Lambda(a\circ\phi_-)\).
Cauchy--Schwarz for this positive semidefinite form gives
\begin{equation}
 |\Lambda(a)|^2
 \le\Lambda(a\circ\phi_+)\Lambda(a\circ\phi_-),
 \qquad
 \Lambda(a\circ\phi_\pm)\ge0.
 \label{eq:folding-cauchy-schwarz}
\end{equation}
This remains valid when the form has a nontrivial null space.
For a constant assignment, \(F=G\); hence \(Z_{\mathcal P}(g)\ge0\).

We next specify a finite sequence of folds whose composition is constant.
For a word \(w\) in block \(b\), write \(p_i=w^{-1}(i)\) for the
position of symbol \(i\).  Let \(\phi_{b,i}\) be the fold onto
\(p_i<p_{i+1}\), where \(1\le i<h\).  Swapping symbols \(i,i+1\)
on the left exchanges precisely the two entries \(p_i,p_{i+1}\).
Thus this fold compares these adjacent entries and puts them in
increasing order.

In each block use the fixed sequence of comparisons
\[
 (1,2),(2,3),\ldots,(h-1,h),
\]
repeated \(h-1\) times.  One sweep moves the largest entry to the
last position.  Once the last \(k\) entries are the largest \(k\)
entries in increasing order, the next sweep leaves them there and puts
the largest remaining entry immediately before them.  Induction shows
that these sweeps sort every position array to \((1,\ldots,h)\).
Concatenating the sequences for all \(\ell\) blocks gives fold maps
\(\phi_1,\ldots,\phi_T\), with \(T=\ell(h-1)^2\), such that
\begin{equation}
 \phi_T\circ\cdots\circ\phi_1(v)=v_*
 \quad\text{for every }v\in V,
 \label{eq:constant-fold-composition}
\end{equation}
where \(v_*\) has the identity word in every block.

To use this sequence, fix any finite nonempty list of real functions
\(f_1,\ldots,f_k\) and set
\[
 M=\max_{a:V\to[k]}
   \left|\mathcal Z_{\mathcal P}((f_{a(v)})_{v\in V})\right|.
\]
The maximum ranges over all assignments, with repetition, and exists
because there are only \(k^d\) assignments.  If \(M>0\), choose a
maximizer.  Each of its folded assignments still takes values in the
same finite list, so each folded integral is at most \(M\) in absolute
value.  Equation~\eqref{eq:folding-cauchy-schwarz} says that those two
integrals are nonnegative and their product is at least \(M^2\).
They must therefore both equal \(M\).  In particular, either fold of
any maximizing assignment is again maximizing.

Apply the folds in reverse sorting order,
\(\phi_T,\phi_{T-1},\ldots,\phi_1\).  Pullbacks compose on the
right, so the resulting assignment is
\(a\circ\phi_T\circ\cdots\circ\phi_1\), which is constant by
\eqref{eq:constant-fold-composition}.  Every intermediate integral is
\(M\), and the last is one of the diagonal values \(Z_{\mathcal P}(f_i)\).
Conversely, every constant assignment was included in the maximum.
It follows that
\begin{equation}
 M=\max_{1\le i\le k}Z_{\mathcal P}(f_i).
 \label{eq:finite-folding-maximum}
\end{equation}
When \(M=0\), all these diagonal values are zero, so the same identity
holds.  This is a finite argument and requires no limiting sequence of
folds. The maximizing-assignment argument is analogous to the proof of
the abstract chessboard estimate in
\cite[Theorem~4.1]{FILS1978ReflectionI}.

Finally, for \(\epsilon>0\), apply
\eqref{eq:finite-folding-maximum} to the list
\[
 f_v=\frac{g_v}{(Z_{\mathcal P}(g_v)+\epsilon)^{1/d}},
 \qquad v\in V.
\]
Every denominator is positive and \(Z_{\mathcal P}(f_v)\le1\).
Multilinearity consequently gives
\[
 \left|\mathcal Z_{\mathcal P}((g_v)_{v\in V})\right|
 \le\prod_{v\in V}(Z_{\mathcal P}(g_v)+\epsilon)^{1/d}.
\]
Letting \(\epsilon\downarrow0\) proves
\eqref{eq:signed-holder}, including zero diagonal values.
\end{proof}

\subsection{The seminorm and its symmetries}

Proposition~\ref{prop:signed-holder} allows us to define, for real \(g\),
\[
 \|g\|_{\mathcal P}=Z_{\mathcal P}(g)^{1/d}.
\]
For \(B\subseteq\mathcal P\), let \(\mathbb E_Bg\) denote
uniform averaging over the coordinates in \(B\), holding the other
coordinates fixed.  We may regard this either as a function on
\(X_{\mathcal P}\) independent of \(B\), or as a function on
\(X_{\mathcal P\setminus B}\).

\begin{lemma}[Seminorm properties]
\label{lem:functional-properties}
The map \(g\mapsto\|g\|_{\mathcal P}\) is a seminorm on the real
functions on \(X_{\mathcal P}\), and
\begin{equation}
 |\mathbb E g|\le\|g\|_{\mathcal P}.
 \label{eq:functional-mean}
\end{equation}
For \(a_p\in\mathbb F_p\) and \(r_p\in Q_p^*\), let
\(A(x)_p=r_px_p+a_p\).  Then
\begin{equation}
 \|g\circ A\|_{\mathcal P}=\|g\|_{\mathcal P}.
 \label{eq:functional-affine}
\end{equation}
For every \(B\subseteq\mathcal P\),
\begin{equation}
 \|\mathbb E_Bg\|_{\mathcal P\setminus B}
   =\|\mathbb E_Bg\|_{\mathcal P}
   \le\|g\|_{\mathcal P}.
 \label{eq:functional-projection}
\end{equation}
\end{lemma}

\begin{proof}
Nonnegativity is already proved.  Since \(d\) is even, homogeneity of
the diagonal form gives \(\|cg\|_{\mathcal P}=|c|\|g\|_{\mathcal P}\)
for every real \(c\).  For real functions \(g_1,g_2\), expand
\(Z_{\mathcal P}(g_1+g_2)\) multilinearly and apply
\eqref{eq:signed-holder} to each term.  With
\(a=\|g_1\|_{\mathcal P}\) and \(b=\|g_2\|_{\mathcal P}\), this gives
\[
 0\le Z_{\mathcal P}(g_1+g_2)
 \le\sum_{S\subseteq V}a^{|S|}b^{d-|S|}=(a+b)^d.
\]
Taking the nonnegative \(d\)-th root proves the triangle inequality,
also when one of the seminorms vanishes.

The probability normalization gives \(Z_{\mathcal P}(1)=1\).
Insert \(g\) in one slot and \(1\) in all others in
\eqref{eq:signed-holder}.  The uniform single-slot marginal identifies
the left side with \(|\mathbb E g|\), proving
\eqref{eq:functional-mean}.  The simultaneous affine invariance of each
\(\mu_p\) proves \eqref{eq:functional-affine}; its parameters may be
chosen independently at different primes.

Let \(G_B\) be the additive subgroup of \(X_{\mathcal P}\) consisting
of vectors supported on \(B\).  Then
\[
 (\mathbb E_Bg)(x)=\frac1{|G_B|}\sum_{t\in G_B}g(x+t).
\]
The triangle inequality and translation invariance show that this average
has seminorm at most \(\|g\|_{\mathcal P}\).  Since the averaged
function is independent of \(B\), integration over the prime laws in
\(B\), each of mass one, gives
\[
 Z_{\mathcal P}(\mathbb E_Bg)
   =Z_{\mathcal P\setminus B}(\mathbb E_Bg).
\]
This proves \eqref{eq:functional-projection}.  All statements include
the empty prime set, for which the seminorm is absolute value.
\end{proof}

\subsection{Conductor decay}

The characters of \(X_{\mathcal P}\) are indexed by
\(\xi\in\mathbb Q/\mathbb Z\) whose reduced denominator \(q(\xi)\)
divides \(\prod_{p\in\mathcal P}p\), with \(q(0)=1\).  We write
such a character as \(e(z\xi)\): it means \(e(n\xi)\) for any integer
\(n\) satisfying \(n\equiv z_p\pmod p\) for all \(p\in\mathcal P\).
This is well defined by the Chinese remainder theorem.  The Fourier
coefficient of \(g\) is
\[
 \widehat g(\xi)=\mathbb E_{z\in X_{\mathcal P}}g(z)e(-z\xi).
\]
For \(U\subseteq\mathcal P\), put \(q_U=\prod_{p\in U}p\),
including \(q_{\varnothing}=1\).  The exact prime support of a
character is the set of primes at which it is nontrivial; its denominator
is exactly the product of these primes.

Recall the exceptional mass and the normalized nonexceptional law,
\[
 \eta_p=\varepsilon_p(\mathbb F_p^V),
 \qquad
 \mu_p^\circ=\frac{\mu_p-\varepsilon_p}{1-\eta_p}.
\]
Proposition~\ref{prop:tuple-laws} gives \(\eta_p\le p^{-4}<1\).

\begin{lemma}[Conductor decay]
\label{lem:conductor-decay}
Let \(\mathcal P\) be a finite set of primes at least \(P\), and
at each prime choose \(\rho_p\) to be either \(\mu_p\) or
\(\mu_p^\circ\).  Let \(R\ge1\), and let
\(g_w:X_{\mathcal P}\to\mathbb C\), \(w\in V\), be functions.
Suppose that for one slot \(v\),
\(\widehat g_v(\xi)=0\) whenever \(q(\xi)\le R\).  Then
\begin{equation}
 \left|\int\prod_{w\in V}g_w(z_w)\,
       d\!\bigotimes_{p\in\mathcal P}\rho_p(z)\right|
 \le R^{-\gamma}\|g_v\|_2
        \prod_{w\ne v}\|g_w\|_{2(d-1)}.
 \label{eq:conductor-decay}
\end{equation}
All norms on the right use uniform probability measure on
\(X_{\mathcal P}\).  The constant in the inequality is one.
\end{lemma}

\begin{proof}
Fix the distinguished slot \(v\).  For each \(p\), let
\(Y_p=\mathbb F_p^{V\setminus\{v\}}\), let \(\lambda_p\) be
the marginal law of the other labels under \(\rho_p\), and define
\[
 T_pf(y)=\mathbb E_{\rho_p}[f(z_v)\mid (z_w)_{w\ne v}=y].
\]
This is an operator from \(L^2(\mathbb F_p,\mathrm{unif})\) to
\(L^2(Y_p,\lambda_p)\); its values on marginal-null points are
irrelevant.  It preserves constants, and the tower property gives
\begin{equation}
 T_p1=1,
 \qquad
 \int T_pf\,d\lambda_p=\mathbb E_{x\in\mathbb F_p}f(x).
 \label{eq:conditional-centering}
\end{equation}
For uniform-mean-zero \(f\), Lemma~\ref{lem:conditional-mixing}
gives
\begin{equation}
 \|T_pf\|_{L^2(\lambda_p)}
   \le p^{-\gamma}\|f\|_{L^2(\mathrm{unif})}.
 \label{eq:local-conductor-contraction}
\end{equation}
Let \(H_p^0\) be the mean-zero subspace of
\(L^2(\mathbb F_p,\mathrm{unif})\), and let \(K_p^0\) be the
mean-zero subspace of \(L^2(Y_p,\lambda_p)\).  Thus
\[
 L^2(\mathbb F_p,\mathrm{unif})=\mathbb C1\mathbin{\oplus}H_p^0,
 \qquad
 L^2(Y_p,\lambda_p)=\mathbb C1\mathbin{\oplus}K_p^0,
\]
with orthogonal sums.  Equation~\eqref{eq:conditional-centering}
says that \(T_pH_p^0\subseteq K_p^0\).  Notice that centering in
the codomain refers to \(\lambda_p\), which need not be uniform.

For \(U\subseteq\mathcal P\), let \(H_U\) be the tensor subspace
with factor \(H_p^0\) when \(p\in U\) and factor \(\mathbb C1\)
otherwise; define \(K_U\) in the same way using \(K_p^0\).
The \(H_U\) form an orthogonal direct sum of
\(L^2(X_{\mathcal P},\mathrm{unif})\).  Distinct \(K_U,K_W\)
are also orthogonal: at a prime in the symmetric difference of \(U,W\),
one tensor factor is constant and the other has mean zero.  Inner products
of pure tensors therefore vanish, and linearity gives the assertion for
the whole subspaces.

Under the product law \(\bigotimes_p\rho_p\), the marginal of all
the other labels is \(\lambda=\bigotimes_p\lambda_p\).  Conditional
independence across primes shows that the operator from the distinguished
slot to all other labels is
\[
 T=\bigotimes_{p\in\mathcal P}T_p.
\]
For completeness, on an input \(\prod_pf_p(z_{v,p})\), the conditional
expectation is \(\prod_pT_pf_p\), since the conditional law is a
product over primes.  Such pure tensors span the finite-dimensional
domain, proving the operator identity.  Consequently
\begin{equation}
 T(H_U)\subseteq K_U,
 \qquad
 \|T|_{H_U}\|\le\prod_{p\in U}p^{-\gamma}=q_U^{-\gamma}.
 \label{eq:exact-support-contraction}
\end{equation}
The norm bound follows by applying
\eqref{eq:local-conductor-contraction} one tensor factor at a time;
on the constant factors the norm is one.

Nontrivial characters at \(p\) form an orthonormal basis for \(H_p^0\).
Thus \(H_U\) consists precisely of Fourier sums with exact denominator
\(q_U\).  Decompose the distinguished input as
\(g_v=\sum_{q_U>R}g_{v,U}\), with \(g_{v,U}\in H_U\).  The
images of distinct summands are orthogonal by
\eqref{eq:exact-support-contraction} and the orthogonality of the
\(K_U\).  Hence
\begin{align*}
 \|Tg_v\|_{L^2(\lambda)}^2
 &=\sum_{q_U>R}\|Tg_{v,U}\|_{L^2(\lambda)}^2\\
 &\le\sum_{q_U>R}q_U^{-2\gamma}\|g_{v,U}\|_2^2
 \le R^{-2\gamma}\|g_v\|_2^2.
\end{align*}
Only distinct exact supports need have orthogonal images; no such
assertion is made for different characters within a single support.

Finally condition the multilinear integral on all labels except those
in slot \(v\), and apply Cauchy--Schwarz.  Its absolute value is at most
\[
 R^{-\gamma}\|g_v\|_2
 \left\|\prod_{w\ne v}g_w(z_w)\right\|_{L^2(\lambda)}.
\]
H\"older's inequality with \(d-1\) factors and the uniform marginal
in each remaining slot give
\[
 \left\|\prod_{w\ne v}g_w(z_w)\right\|_{L^2(\lambda)}
 \le\prod_{w\ne v}
       \left(\mathbb E_{\lambda}|g_w(z_w)|^{2(d-1)}\right)^{1/(2(d-1))}
 =\prod_{w\ne v}\|g_w\|_{2(d-1)}.
\]
This proves \eqref{eq:conductor-decay}.  If \(\mathcal P\) is empty,
the support hypothesis forces \(g_v=0\), so the assertion holds there
too.  No reflection positivity of the laws \(\rho_p\) was used.
\end{proof}
